%% LyX 2.4.4 created this file.  For more info, see https://www.lyx.org/.
%% Do not edit unless you really know what you are doing.
\documentclass[english]{article}
\usepackage[T1]{fontenc}
\usepackage[latin9]{inputenc}
\usepackage{enumitem}
\usepackage{amsmath}
\usepackage{amsthm}
\usepackage{amssymb}
\usepackage{geometry}
\geometry{verbose,tmargin=1in,bmargin=1in,lmargin=1in,rmargin=1in}
\usepackage{wasysym}

\makeatletter
%%%%%%%%%%%%%%%%%%%%%%%%%%%%%% Textclass specific LaTeX commands.
\numberwithin{equation}{section}
\numberwithin{figure}{section}
\newlength{\lyxlabelwidth}      % auxiliary length 

%%%%%%%%%%%%%%%%%%%%%%%%%%%%%% User specified LaTeX commands.
\usepackage{fancyhdr}
\usepackage{lastpage}
\pagestyle{fancy}
\usepackage{enumitem}
\usepackage{ifthen}
\everymath=\expandafter{\the\everymath\displaystyle}
%\DeclareMathSizes{10}{10}{10}{10}
\usepackage{multicol}

\usepackage{tikz}
\usetikzlibrary{patterns}
\usetikzlibrary{plotmarks}
\usepackage{pgfplots}
\pgfplotsset{compat=newest}

\usepackage{calc}
\usepackage[nomessages]{fp}

\usepackage{datetime}
% Added by lyx2lyx
\setlength{\parskip}{\smallskipamount}
\setlength{\parindent}{0pt}

\makeatother

\usepackage{babel}
\begin{document}
\renewcommand{\author}{Dr.\,James Ashe}

\newcommand{\classNum}{335}

\newcommand{\classSec}{A}

\newcommand{\nameType}{Name: \hspace{-4cm}}

\newcommand{\examType}{Homework 2}

\newcommand{\Date}{\formatdate{28}{1}{2014}} %%\AdvanceDate[12] \today

%\newdate{my_date}{\day}{\month}{\year}%   
%\AdvanceDate[-5] 
%\displaydate{my_date}

%%First page's header%%%%%%%%%%%%%%%%%%%%%%
\fancypagestyle{firststyle} %%header settings for first pagr
{
	\fancyhead[L]{MTH \classNum{}(\classSec{}) \author{}\\
	\textbf{\examType{}} \Date{}}
	\fancyhead[C]{\textbf{\nameType}}
	\setlength{\headsep}{14pt}
}
%%%%%%%%%%%%%%%%%%%%%%%%%%%%%%%%%%%%%%%%%%

%%Next page's header%%%%%%%%%%%%%%%%%%%%%%%
%\setlist{nolistsep} %eliminates space between list items
\lhead{MTH \classNum{}(\classSec{})} 
\chead{\textbf{\nameType}} 
\cfoot{\thepage{}~of~\pageref{LastPage}} 
\setlength{\headsep}{5pt} 
%%%%%%%%%%%%%%%%%%%%%%%%%%%%%%%%%%%%%%%%%%%

%%Various settings; see the preamble for other settings%%%%%
%%\setenumerate[0]{leftmargin=12pt,itemindent=0pt,labelwidth=0pt}
\renewcommand{\labelenumi}{\textbf{\arabic{enumi}.}}
\renewcommand{\labelenumii}{\textbf{\alph{enumii}.}}
\thispagestyle{firststyle}
%%%%%%%%%%%%%%%%%%%%%%%%%%%%%%%%%%%%%%%%%%

\textbf{Homework}
\begin{enumerate}
\item Prove or disprove the following.

\begin{enumerate}[resume]
\item A constant function is not one-to-one if the domain has more than
one element.
\item The identity function is one-to-one.
\item The identity function is onto.\vfill
\end{enumerate}
\item Let $g:A\rightarrow B$ and $f:B\rightarrow C$. Prove that if $f\circ g:A\rightarrow C$
is one-to-one and $g$ is onto then $f$ is also one-to-one. \vfill
\item Suppose that $f:A\rightarrow B$ is a one-to-one function. Prove or
disprove that if $f^{-1}:f(B)\rightarrow A$ is a function then $f^{-1}$
is also one-to-one.\vfill
\end{enumerate}
\begin{enumerate}[resume]
\item Let $f:A\rightarrow B$ be a one-to-one function. Show that $\left|A\right|\leq\left|B\right|$.\vfill
\item Let $f:A\rightarrow B$ and $g:B\rightarrow C$ be functions. Prove
that $g\circ f:A\rightarrow C$ is a function.\vfill
\item (This problem was dropped) Let $f:A\rightarrow A$, $g:A\rightarrow A$,
and $h:A\rightarrow A$; prove that $\left(h\circ g\right)\circ f=h\circ\left(g\circ f\right)$.\vfill
\item Prove that if $f:A\rightarrow B$ is onto and $g:B\rightarrow C$
is onto, then $\left(g\circ f\right):A\rightarrow C$ is onto. \vfill
\item Each of the following statements, $P(x,y)$, define a relation $R$
in the natural numbers, $\mathbb{N}$, i.e., $R=(\mathbb{N},\mbox{\ensuremath{\mathbb{N}}},P(x,y))$.
Show whether each of the following relations is reflexive. 

\begin{enumerate}
\item $P(x,y)=$``$x$ is less than or equal to $y$''\vfill
\end{enumerate}
\begin{enumerate}[resume]
\item $P(x,y)=$``$x$ divides $y$''\vfill
\item $P(x,y)=$``$x+y=10$''\vfill
\item $P(x,y)=$ ``$x$ and $y$ are relatively prime.''\vfill
\end{enumerate}
\end{enumerate}
\textbf{Presentation problems.}
\begin{enumerate}
\item Let the function $f:A\rightarrow B$ be 1-1 and onto, i.e., $f^{-1}$
exists and $\left|A\right|=\left|B\right|$. Prove that $\left(f^{-1}\circ f\right):A\rightarrow A$
is the identity function on $A$ and that $\left(f\circ f^{-1}\right):B\rightarrow B$
is the identity function on $B$.
\item Let $f:A\rightarrow B$ be 1-1 and onto. Prove that $g=f^{-1}$, if
$\left(g\circ f\right):A\rightarrow A$ is the identity function on
$A$ and $\left(f\circ g\right):B\rightarrow B$ is the identity function
on $B$. 
\item The \emph{product set }of sets $A$ and $B$ is the set of all ordered
pairs $(a,b)$ where $a\in A$ and $b\in B$; it is denoted by $A\times B$.
Prove that $A\times(B\cap C)=(A\times B)\cap(A\times C)$. 
\item Prove: Let $R$ and $R'$ be symmetric relations in a set $A$; then
$R\cap R'$ is a symmetric relation in $A$.
\item Prove that ${\displaystyle \sum_{k=1}^{n}2k=n^{2}+n}$
\item The Fibonacci numbers are :$1,1,2,3,5,8,13,21,34,\dots$. In general,
the Fibonacci numbers are defined by $f_{1}=1$, $f_{2}=1$, and $f_{n+2}=f_{n+1}+f_{n}$
for $n=1,2,3,\dots$. Prove that the $n^{\mbox{th}}$ Fibonacci number
$f_{n}<2^{n}$.
\item Prove that there infinitely many prime numbers. 
\item [$\smiley$\textbf{.}]Two sets are said to be \emph{equal }if there
exists a bijective function between them; we write $\left|A\right|=\left|B\right|$.
$B$ is said to have \emph{strictly greater cardinality} than $B$
if there is a one-to-one function from $A$ into $B$, but there does
not exist a bijective function from $A$ into $B$; we write$\left|A\right|<\left|B\right|$.
Prove that $\left|\mathbb{N}\right|<\left|\mathbb{R}\right|$.
\end{enumerate}
\vfill
\end{document}
